\documentclass[../../main.tex]{subfiles}
\graphicspath{{\subfix{../../image/}}}

\begin{document}

\section{偏微分方程的研究策略}\label{section:ch01-strategies}

如 \S1.1 所说明，我们的目标是发现求解各类偏微分方程的方法；但鉴于 \S1.2 中罗列的众多五花八门的例子，现在应当清楚：这绝非易事。事实上，``求解''一个给定 PDE 究竟意味着什么，这个问题本身就可能是微妙的，它在很大程度上取决于所讨论问题的特定结构。

\subsection{适定问题与经典解}\label{subsection:ch01-wellposed}

一个适定问题的直观概念，抓住了``求解一个 PDE''所应具备的许多令人希望的特征。我们称一个偏微分方程的给定问题是\textbf{适定的}，如果
\begin{enumerate}[(i)]
\item 该问题确实有解；
\item 此解唯一；且
\item 解连续地依赖于问题所给的数据。
\end{enumerate}
最后一个条件对于来自物理应用的问题尤为重要：当规定问题的条件有少许改变时，我们希望（唯一）解只改变少许。（另一方面，对许多问题并不指望有唯一性；在这些情形中，主要的数学任务是对解进行分类与刻画。）

显然，若能以满足 (i)--(iii) 的方式``求解''PDE，那当然再好不过。但请注意，我们仍未仔细地定义所谓``解''究竟指什么。例如，我们是否应当要求一个``解'' $u$ 是实解析的，或者至少是无穷次可微的？这或许很理想，但也许我们要求得太多了。或许更明智的做法是：要求一个 $k$ 阶 PDE 的解至少 $k$ 次连续可微。这样，至少 PDE 陈述中出现的所有导数都存在且连续，尽管某些更高阶导数可能不存在。我们非正式地把具有这种光滑性的解称为该 PDE 的\textbf{经典解}：这当然是最显而易见的解的概念。

因此，所谓在经典意义下求解一个偏微分方程，就是在可能时写出一个满足上述 (i)--(iii) 的经典解的公式，或至少证明这样一个解存在，并推导出它的各种性质。

\subsection{弱解与正则性}\label{subsection:ch01-weak}

但我们能否做到这一点呢？答案是：某几类特定的偏微分方程（例如Laplace方程）可以在经典意义下求解，但其他许多方程（即便不是大多数）则不能。例如考虑标量守恒律
\begin{align*}
u_t+F(u)_x=0.
\end{align*}
我们将在 \S3.4 中看到，这个 PDE 支配着涉及流体力学的各种一维现象，特别是刻画了激波的形成与传播。激波是解 $u$ 的一条间断曲线；因此，若想研究守恒律并还原其背后的物理，就必须允许 $u$ 不必连续可微、甚至不必连续。一般而言，正如我们将看到的，守恒律没有经典解，但只要允许适当定义的广义解或弱解，它就是适定的。

这一切都说明：方程特定的结构可能迫使我们放弃对光滑经典解的寻求。我们必须在仍希望达到适定性条件 (i)--(iii) 的同时，去考察更广泛的候选解类。而事实上，即使对那些最终可在经典意义下求解的 PDE，起初先去寻找某种适当的弱解也往往是最方便的。

关键在于：倘若从一开始就要求解非常正则，比如说 $k$ 次连续可微，那么我们通常会在寻找解时陷入真正的困境，因为证明就必然包含可能十分繁复的论证，以说明我们所构造的函数确实足够光滑。一个远为合理的策略是把存在性问题与光滑性（或正则性）问题分开处理。其想法是：对给定的 PDE 定义一个足够宽泛的弱解概念，并期望由于我们对弱解的光滑性要求并不太高，从而较容易建立它的存在性、唯一性以及对所给数据的连续依赖性。因此，重复一遍，以在某个适当的弱解或广义解类中证明适定性为目标，往往是明智的。

如上所述，对某些偏微分方程而言，这已是所能做到的最好结果。对另一些方程，我们则可以希望弱解最终终究足够光滑，从而可成为经典解。这就引出弱解的\textbf{正则性}问题。正如我们将看到的，弱解的存在性常常依赖于相当简单的估计加上泛函分析的思想，而弱解的正则性（当其成立时）则通常依赖于许多精巧的微积分估计。

在此我要明确指出：一旦过了第一部分（\cref{chapter:ch02,chapter:ch03,chapter:ch04}），我们的工作将主要致力于从数学上证明各类偏微分方程解的存在性，而不那么着重于推导这些解的公式。这看起来像是徒劳或误入歧途，但实际上数学家就像神学家：我们把存在性视为所研究对象的首要属性。但与神学家不同，我们不必总是仅靠信仰行事。

\subsection{典型困难}\label{subsection:ch01-difficulties}

下面是一些笼统但普遍的原则，记住它们或许有益：
\begin{enumerate}[(i)]
\item 非线性方程比线性方程更难；而且，非线性性越是影响到高阶导数，PDE 就越难。
\item 高阶 PDE 比低阶 PDE 更难。
\item 方程组比单个方程更难。
\item 含多个自变量的偏微分方程比含少数自变量的 PDE 更难。
\item 对大多数偏微分方程而言，不可能写出解的显式公式。
\end{enumerate}
这些论断无一没有重要的例外。

\end{document}